% Part: lambda-calculus
% Chapter: introduction
% Section: church-rosser

\documentclass[../../../include/open-logic-section]{subfiles}

\begin{document}

\olfileid[pt]{lam}{int}{cr}
\olsection{A propriedade de Church--Rosser}

\begin{thm}
\ollabel{thm:church-rosser}
Sejam $M$, $N_1$ e $N_2$ termos tais que $M \red N_1$ e $M \red N_2$.
Então existe um termo $P$ tal que $N_1 \red P$ e $N_2 \red P$.
\end{thm}

\begin{cor}
Suponha que $M$ possa ser reduzido à forma normal. Então essa forma normal é
única.
\end{cor}

\begin{proof}
Se $M \red N_1$ e $M \red N_2$, pelo teorema anterior existe um termo~$P$
tal que $N_1$ e $N_2$ se reduzem ambos a~$P$. Se $N_1$ e~$N_2$ estão ambos
em forma normal, isso só pode acontecer se $N_1 \ident P \ident N_2$.
\end{proof}

Por fim, diremos que dois termos $M$ e~$N$ são \emph{$\beta$-equivalentes},
ou simplesmente \emph{equivalentes}, se eles se reduzem a um termo comum;
em outras palavras, se existe algum $P$ tal que $M \red P$ e $N \red P$.
Escrevemos isso como $M \equal[\beta] N$. Usando
\olref{thm:church-rosser}, você pode verificar que $\equal[\beta]$ é uma
relação de equivalência com a propriedade adicional de que, para quaisquer
$M$ e~$N$, se $M \red N$ ou $N \red M$, então $M \equal[\beta] N$. (Na
verdade, pode-se mostrar que $\equal[\beta]$ é a \emph{menor} relação de
equivalência com essa propriedade.)

\end{document}


